\documentclass[letterpaper]{article} 
\usepackage[utf8]{inputenc}
\linespread{0.85}
\usepackage[T1]{fontenc}
\usepackage{amsmath}
\usepackage{amsfonts}
\usepackage{amssymb}
\usepackage{array}
\usepackage{fontspec}
\usepackage{booktabs}
\usepackage{hyperref}
\usepackage[version=4]{mhchem}
\usepackage{stmaryrd}
\usepackage[dvipsnames]{xcolor}
\colorlet{LightRubineRed}{RubineRed!70}
\colorlet{Mycolor1}{green!10!orange}
\definecolor{Mycolor2}{HTML}{00F9DE}
\usepackage{graphicx}
\usepackage{amsmath}
\usepackage{graphicx}
\usepackage{capt-of}
\usepackage{lipsum}
\usepackage{fancyvrb}
\usepackage{tabularx}
\usepackage{listings}
\usepackage[export]{adjustbox}
\graphicspath{ {./images/} }
\usepackage[utf8]{inputenc}
\usepackage[english]{babel}
\usepackage{float}
\usepackage{lipsum}
\usepackage{graphicx}
\usepackage{float}
\usepackage[margin=0.7in]{geometry}
\usepackage{amsmath}
\usepackage{graphicx}
\usepackage{capt-of}
\usepackage{tcolorbox}
\usepackage{lipsum}
\usepackage{graphicx}
\usepackage{float}
\usepackage{listings}
\definecolor{deepred}{RGB}{139,0,0}
\usepackage{hyperref} 
\usepackage{xcolor} % For custom colors
\lstset{
	language=Python,                % Choose the language (e.g., Python, C, R)
	basicstyle=\ttfamily\small, % Font size and type
	keywordstyle=\color{blue},  % Keywords color
	commentstyle=\color{gray},  % Comments color
	stringstyle=\color{red},    % String color
	numbers=left,               % Line numbers
	numberstyle=\tiny\color{gray}, % Line number style
	stepnumber=1,               % Numbering step
	breaklines=true,            % Auto line break
	backgroundcolor=\color{black!5}, % Light gray background
	frame=single,               % Frame around the code
}
\usepackage{float}
\usepackage[]{amsthm} %lets us use \begin{proof}
	\usepackage[]{amssymb} %gives us the character \varnothing

	\title{Homework 7, MATH 4155}
	\author{Zongyi Liu}
	\date{Tue, Feb 24, 2026}
	\begin{document}
		\maketitle
		\section{Question 1}
		
		\begin{itemize}
			\item[(i)] Suppose $\{S_n\}_{n\in\mathbb{N}_0}$ is the simple random walk started at $S_0=0$ (cf.\ Chapter 9 of Notes). What is the probability
			\[
			\mathbb{P}\bigl(S_1\ge 0,\cdots,S_{N-1}\ge 0 \mid S_N=x\bigr)
			\]
			that the random walk never took a negative value in between, given that it landed at the integer site $x>0$ on day $N$?
			
			\item[(ii)] Given that a game of ping-pong finished 21--18, what is the probability that the eventual victor never fell behind throughout the entire match? State very precisely your assumptions.
		\end{itemize}
		
		\textbf{Answer}
		
		\clearpage
		
		\section{Question 2}
		
		For the simple random walk $\left\{S_{n}\right\}_{n \in \mathbb{N}_{0}}$ started at $S_{0}=0$, consider the time
		
		
		\begin{equation*}
			T_{b}:=\inf \left\{n \in \mathbb{N} \mid S_{n}=b\right\} \tag{0.1}
		\end{equation*}
		
		
		of first visit (respectively, of first return) to the level $b \in \mathbb{N}$ (resp., $b=0$ ). As always, we take $\inf \emptyset=\infty$.\\
		(i) For any given $k \in \mathbb{N}$, show that we have
		
		$$
		\mathbb{P}_{0}\left(T_{k}=k+2 n\right)=\frac{k}{k+2 n}\binom{k+2 n}{n} p^{k+n} q^{n}, \quad n \in \mathbb{N}_{0}
		$$
		
		(ii) In particular, for $p=q=1 / 2$ and $k=1$ show
		
		$$
		\mathbb{P}_{0}\left(T_{1}=2 n-1\right)=\frac{1}{2 n-1}\binom{2 n-1}{n}\left(\frac{1}{2}\right)^{2 n-1}=\mathbb{P}_{0}\left(T_{0}=2 n\right), \quad n \in \mathbb{N} .
		$$
		
		\textbf{Answer}
		
		\clearpage
		
		\section{Question 3}
		
		Argue from first principles that convergence in probability implies vague convergence.
		
		\textbf{Answer}
		
		\clearpage
		
		\section{Question 4}
		
		{\color{deepred}\fontspec{Georgia}  Markov-Cantelli Law of Large Numbers}
		
		 Consider independent random variables $X_{1}, X_{2}, \cdots$ with $\sup _{n \in \mathbb{N}} \mathbb{E}\left(X_{n}^{4}\right)<\infty$. Show that
		
		$$
		\lim _{n \rightarrow \infty} \frac{1}{n} \sum_{j=1}^{n}\left(X_{j}-\mathbb{E}\left(X_{j}\right)\right)=0, \quad \text { holds a.e. }
		$$
		
		(\emph{Hint}: Use Čebyšev, Borel-Cantelli and Theorem 10.1 in the Notes. Start with the case of identically distributed random variables. Note also that the full strength of independence is not necessary.)
		
		
		\textbf{Answer}
		
		\clearpage
		
		\section{Question 5}
		
		Let $X_{1}, X_{2}, \cdots$ be independent random variables with $\mathbb{P}\left(X_{n}=1\right)=p_{n} \in (0,1)$ and $\mathbb{P}\left(X_{n}=0\right)=1-p_{n}$. Show that $\lim _{n \rightarrow \infty} X_{n}=0$ holds
		
		\begin{itemize}
			\item in probability, if and only if $\lim _{n \rightarrow \infty} p_{n}=0$;
			\item in $\mathbb{L}^{r}$ for $r>0$, if and only if $\lim _{n \rightarrow \infty} p_{n}=0$;
			\item a.e., if and only if $\sum_{n \in \mathbb{N}} p_{n}<\infty$.
		\end{itemize}
		
		\textbf{Answer}
		
		\clearpage
		
		\section{Question 6}
		
		{\color{deepred}\fontspec{Georgia} Walsh 5.10} 
		
		Let $(X_n)$ be random variables and let $S_n=X_1+\cdots+X_n$. Suppose that $X_n\to 0$ a.e. Show that $S_n/n\to 0$. Show that this is no longer true for convergence in probability: $X_n\to 0$ in probability does not imply $S_n/n\to 0$ in probability, even if the $X_n$ are independent.
		
		(\emph{Hint}: Let $X_n$ equal $0$ and $2^n$ with probabilities $1-1/n$ and $1/n$ respectively.)
		
		\textbf{Answer}
		
		\clearpage
		
		\section{Question 7}
		
		{\color{deepred}\fontspec{Georgia} Walsh 5.23} 
	
	Let $X_1,X_2,\cdots$ be a sequence of i.i.d.\ random variables with $E\{|X_1|\}=\infty$.
	\begin{enumerate}
	\item[(a)] Show that $\displaystyle \limsup_{n\to\infty}\frac{|X_n|}{n}=\infty$ a.s.
	\item[(b)] Deduce that $\displaystyle \limsup_{n\to\infty}\frac{|X_1+\cdots+X_n|}{n}=\infty$.
\end{enumerate}

(\textit{Hint:} Recall the ``layered representation''
\[
\mathbb{E}(Z)=\int_0^\infty \mathbb{P}(Z>u)\,du
\le \sum_{n\in\mathbb{N}_0}\mathbb{P}(Z>n)
\le 1+\mathbb{E}(Z)
\]
for the expectation of a random variable $Z:\Omega\to[0,\infty)$. Use---what else---the (second) \textsc{Borel--Cantelli} Lemma.)

		
		\textbf{Answer}
		
		\clearpage
		
		\section{Question 8}
		
		For the simple random walk starting at the origin $x=0$ (with no barriers), find the probability that a given integer $b\in\mathbb{N}$ will eventually be reached.
		
		\textbf{Answer}
		
		\clearpage
		
		\section{Question 9}
		
		Let $S_{0}, S_{1}, S_{2}, \cdots$ be an unconstrained (``no barriers") simple random walk starting at the origin $S_{0}=0$ (cf. Chapter 9 of Notes). Show that
		
		$$
		\mathbb{P}\left(S_{1} \geq 0, S_{2} \geq 0 \cdots, S_{2 n-1} \geq 0, S_{2 n}=0\right)=\frac{1}{n+1} u_{2 n}
		$$
		
		where
		
		$$
		u_{2 n}:=\frac{(2 n)!}{n!n!}(p q)^{n}
		$$
		
		is the probability of $n$ successes and $n$ failures in $2 n$ Bernoulli trials (coin tosses). Deduce
		
		$$
		\mathbb{P}\left(S_{1} \geq 0, S_{2} \geq 0 \cdots, S_{2 n-1} \geq 0 \mid S_{2 n}=0\right)=\frac{1}{n+1}
		$$
		
		\textbf{Answer}
		
		\clearpage
		
		\section{Question 10}
		
		We place ourselves in the same setting as in Problem \#9, and recall the notation established there (cf. Chapter 9 of Notes). We also recall from class the computation
		
		$$
		\begin{aligned}
			h_{2 n}:=\mathbb{P}\left(S_{1} \neq 0, S_{2} \neq 0 \cdots, S_{2 n-1} \neq 0, S_{2 n}=0\right) & =2 \mathbb{P}\left(S_{1}>0, S_{2}>0 \cdots, S_{2 n-1}>0, S_{2 n}=0\right) \\
			& =\frac{2(p q)^{n}}{n} \frac{(2 n-2)!}{(n-1)!(n-1)!}=\frac{2 p q}{n} u_{2 n-2}
		\end{aligned}
		$$
		
		as well as the Stirling approximation
		
		$$
		k!\sim \sqrt{2 \pi k}(k / e)^{k}, \quad k \in \mathbb{N}
		$$
		
		for the factorial of an integer, in the sense that the ratio of the left- over the right-hand side converges to 1 as $n \rightarrow \infty$.
		
		We denote the time of first return to the origin for the random walk by
		
		$$
		R:=\inf \left\{k \geq 1: S_{k}=0\right\}
		$$
		
		We employ always the convention, that the infimum of the empty set is $+\infty$. In other words, on the event $\{R=\infty\}$ the random walk never returns to the origin.
		
		\begin{itemize}
			\item[(i)] What is the distribution of the random variable $R$?
			
			\item[(ii)] Show that the expected time of first return to the origin is infinite, i.e., $\mathbb{E}(R)=\infty$, in the ``symmetric case'' $p=q=1/2$.
			
			\item[(iii)] Verify that the generating function
			\[
			G(s):=\mathbb{E}(s^{R})
			=\mathbb{E}\bigl(s^{R}\mathbf{1}_{\{R<\infty\}}\bigr)
			=\sum_{n\in\mathbb{N}}\mathbb{P}(R=2n)s^{2n}
			=\sum_{n\in\mathbb{N}} h_{2n}s^{2n},
			\qquad 0<s<1,
			\]
			is given as
			\[
			G(s)=1-\sqrt{\,1-4pqs^{2}\,}.
			\]
			
			\item[(iv)] Recover from the computation in (iii) a result shown in class: The probability of the event $\{R=\infty\}$, that the random walk never returns to the origin, is
			\[
			\mathbb{P}(R=\infty)=|p-q|.
			\]
			Deduce that, in the symmetric case $p=q=1/2$, the random walk returns eventually to the origin with probability one---even though this ``may take an awfully long time'', in the sense that $\mathbb{E}(R)=\infty$ holds as we have seen in (ii).
			
			\item[(v)] Deduce that $\mathbb{E}(R)=\infty$ holds always, not just in the symmetric case $p=q=1/2$.
		\end{itemize}
		
		
		\textbf{Answer}
		
		\clearpage
		
		
	\end{document}
